%!TEX root = forallxyyc.tex

\chapter{符号记法}
\label{app.notation}

\section{其他命名习惯}

\paragraph{真值函项逻辑} 真值函项逻辑还有其他名称。有时它被称为\emph{语句逻辑}，因为它从根本上处理的是语句。有时它被称为\emph{命题逻辑}，基于它从根本上处理的是命题这一观点。我们采用了\emph{真值函项逻辑}这一名称，以强调它只处理语句的真假赋值，并且其所有联结词都是真值函项的。

\paragraph{一阶逻辑} 一阶逻辑也有其他名称。有时它被称为\emph{谓词逻辑}，因为它允许我们将谓词应用于对象。有时它被称为\emph{量化逻辑}，因为它使用了量词。

\paragraph{公式} 有些文献将公式称为\emph{良构公式}。由于“良构公式”这个短语冗长繁琐，便将其缩写为\emph{wff}。这种做法既笨拙也无必要（这类文献并不承认“非良构公式”）。我们坚持使用“公式”这一说法。

在\cref{s:TFLSentences}中，我们定义了真值函项逻辑的\emph{语句}。这些有时也被称为“公式”（或“良构公式”），因为在真值函项逻辑中，与一阶逻辑不同，公式和语句之间没有区分。

\paragraph{赋值} 有些文献将赋值称为\emph{真值指派}或\emph{真值赋值}。

\paragraph{$n$元谓词} 我们选择将谓词称为“一元”、“二元”、“三元”等。其他文献则分别称其为“一元（monadic）”、“二元（dyadic）”、“三元（triadic）”等。还有些文献称其为“一元（unary）”、“二元（binary）”、“三元（ternary）”等。

\paragraph{名称} 在一阶逻辑中，我们用“$a$”、“$b$”、“$c$”等表示名称。有些文献称它们为“常元”。另一些文献则在句法上不区分名称和变元。这些文献仅关注符号是\emph{被约束的}还是\emph{自由的}。

\paragraph{论域} 有些文献将论域描述为“话语论域”或“论域宇宙”。

\section{其他符号}
在形式逻辑史上，不同时期、不同作者使用了不同的符号。作者常常不得不使用其印刷商能够排版的记法。本附录介绍一些常见符号，以便你在其他文章或书籍中遇到时能够识别。

\paragraph{否定} 两个常用的符号是\emph{锄头形}符号“$\neg$”和\emph{波浪线}符号“${\sim}$”。在一些更高级的形式系统中，有必要区分两种否定；这种区分有时通过同时使用“$\neg$”和“${\sim}$”来表示。较老的文献有时通过在被否定的公式上方画线（例如“$\overline{A \eand B}$”）或使用减号“$-$”来表示否定。许多文献用“$x \neq y$”来缩写“$\enot x = y$”。

\paragraph{析取} 符号“$\vee$”通常用于表示相容析取。在一些所谓代数传统的文献中，析取被写作加号“$+$”。在许多编程语言中，使用竖线“\verb+|+” （或双竖线“\verb+||+”）。（这容易与\cref{c:FunctionalCompleteness}中介绍的谢弗竖线混淆，后者通常写作“$\mid$”。）

\paragraph{合取}
合取通常用\emph{与号}“{\&}”表示。与号是拉丁词“et”（意为“和”）的一种装饰性写法。（其词源在某些字体中仍有体现，尤其是在斜体字体中；因此一个斜体的与号可能显示为“%
\iflatexml
\<span class="ltx\_text" style="font-family:Palatino; font-style:italic;">\&\</span>
\else
\emph{\&}
\fi
”。）这个符号在自然英语书写中很常见（例如“Smith \& Sons”），因此，尽管它是很自然的选择，但许多逻辑学家为了避免对象语言与元语言混淆而使用不同的符号：作为形式系统内的符号，与号并非英语单词“\&”。如今最常见的选择是“$\wedge$”，它与用于析取的符号相呼应。有时也使用一个点号“{\scriptsize\textbullet}”。在一些更老的文献中，合取根本没有专门的符号；“$A$ 且 $B$”直接被写作“$AB$”。

\paragraph{实质条件} 实质条件有两个常用符号：\emph{箭头}“$\rightarrow$”和\emph{马蹄形}“$\supset$”。

\paragraph{实质双条件} \emph{双头箭头}“$\leftrightarrow$”用于那些用箭头表示实质条件的系统。使用马蹄形表示条件的系统则通常用\emph{三线等号}“$\equiv$”来表示双条件。

\paragraph{量词} 全称量词通常被符号化为一个旋转的“A”，存在量词被符号化为一个旋转的“E”。在一些文献中，全称量词没有单独的符号。取而代之的是，变元被直接写在它所约束的公式前面的括号内。例如，他们可能会写“$(x)\atom{P}{x}$”，而我们则会写“$\forall x\,
\atom{P}{x}$”。一些文献也分别用我们合取和析取联结词的大号版本，即“$\bigwedge$”和“$\bigvee$”，来表示全称量词和存在量词。（变元有时被设为这些符号的下标，甚至写在它们正下方。）

\bigskip

这些不同的排印方式总结如下：

\begin{center}
\begin{tabular}{rl}
否定 & $\neg$, ${\sim}$, $-$, $\overline{A}$\\
合取 & $\wedge$, $\&$, {\scriptsize\textbullet}\\
析取 & $\vee$, $+$, \verb+|+, \verb+||+\\
条件 & $\rightarrow$, $\supset$\\
双条件 & $\leftrightarrow$, $\equiv$\\
全称量词 & $\forall x$, $(x)$, $\bigwedge_x$\\
存在量词 & $\exists x$, $\bigvee_x$
\end{tabular}
\end{center}

\section{波兰记法}

本节简要讨论使用波兰记法的语句逻辑。波兰记法是一种记法系统，由波兰逻辑学家扬·武卡谢维奇于1920年代末引入。

小写字母用作语句字母。大写字母 $N$ 用于表示否定。$A$ 用于表示析取，$K$ 用于表示合取，$C$ 用于表示条件，$E$ 用于表示双条件。（“$A$”代表交替，即逻辑析取的另一个名称。“$E$”代表（等价）。）

在波兰记法中，二元联结词被写在它所连接的两个语句\emph{之前}。例如，真值函项逻辑中的语句 $A\eand B$ 在波兰记法中将被写作 $Kab$。

语句 $\enot A\eif B$ 和 $\enot (A\eif B)$ 非常不同：第一个的主逻辑算子是条件，但第二个的主逻辑算子是否定。在真值函项逻辑中，我们通过在第二个语句中给条件句加上括号来表明这一点。在波兰记法中，则完全不需要括号。最左边的联结词总是主逻辑算子。第一个语句可以直接写作 $CNab$，而第二个则写作 $NCab$。

波兰记法的这一特性意味着，只需从右向左依次处理符号，就可以对语句进行求值。例如，若要为 $NKab$ 构造真值表，你需要先考虑赋给 $b$ 和 $a$ 的真值，然后计算它们的合取，最后再将结果否定。而在真值函项逻辑中，确定下一步求值对象的一般规则则远没有这么简单。在真值函项逻辑中，为 $\enot(A\eand B)$ 构造真值表需要先看 $A$ 和 $B$，然后看位于语句中间的合取符号，最后再看语句开头的否定符号。由于波兰记法中的运算顺序可以更机械地规定，其变体被用作许多计算机编程语言的内部结构。